\documentclass[a4paper,10pt]{article}
\usepackage[T1]{polski}
\usepackage[utf8]{inputenc}
\usepackage{xcolor}
\usepackage{titlesec}
\usepackage{graphicx}
\usepackage[inkscapeformat=pdf]{svg}
\usepackage{minted}
\usepackage{listings}
\lstdefinestyle{custom}{
    backgroundcolor=\color{white},
    inputencoding=utf8,
    extendedchars=\true,
    commentstyle=\color{red},
    keywordstyle=\color{blue},
    breaklines=true,
    basicstyle=\ttfamily\footnotesize,
    tabsize=1,
    escapechar='
}
\lstset{style=custom}
\titlespacing*{\section}{0pt}{6pt}{3pt}
\titlespacing*{\subsection}{0pt}{4pt}{2pt}
\hoffset=-3.0cm
\textwidth=18cm
\evensidemargin=0pt
\voffset=-3cm
\textheight=27cm
\setlength{\parindent}{0pt}
\setlength{\parskip}{\medskipamount}
\raggedbottom
\newcommand\BS{\char`\\}
\newcommand\TY{\raise.17ex\hbox{$\scriptstyle\mathtt{\sim}$}}
\usepackage{float}

\author{Matylda Drapich}
\title{Laboratorium sieci komputerowych\\ćw.5\\Konfiguracja serwera w chmurze}
\frenchspacing

\begin{document}
\thispagestyle{empty}
\maketitle

\vspace{-0.5em}
Celem ćwiczenia była skryptowa konfiguracja serwera VPS w chmurze Azure
oraz podłączenie urządzeń do prywatnej sieci ZeroTier ADM.
Przygotowano trzy skrypty w \texttt{/bin/sh}: \texttt{gen-vps}, \texttt{config-vps}
oraz \texttt{config-www}. Serwer działa w regionie PolandCentral
z publicznym adresem \texttt{74.248.150.53} i adresem ADM \texttt{10.80.108.71}.

\section{Skryptowa konfiguracja serwera w chmurze}

\subsection{Generacja VPS -- \texttt{gen-vps}}

Skrypt \texttt{gen-vps} rejestruje zasoby Azure (\texttt{az/register}, \texttt{az/nsg -a}),
tworzy maszynę wirtualną Archlinux (\texttt{az/archlinux -k}) i konfiguruje
mapowanie nazw na stacji roboczej. Adres DNS serwera trafia do pliku
\texttt{\TY/.vps}, a nazwa \texttt{vps} jest mapowana w \texttt{/etc/hosts}
na publiczny adres IP. Dzięki temu możliwe jest logowanie poleceniem
\texttt{ssh vps} bez podawania pełnej nazwy DNS.

\begin{lstlisting}
drapich@s3 ~ % gen-vps
ResourceGroup  PowerState   PublicIpAddress  Fqdns
LabSK          VM running   74.248.150.53    drapich.polandcentral.cloudapp.azure.com
74.248.150.53 vps
drapich@s3 ~ % cat ~/.vps
drapich.polandcentral.cloudapp.azure.com
drapich@s3 ~ % grep vps /etc/hosts
74.248.150.53 vps
\end{lstlisting}

\subsection{Konfiguracja VPS -- \texttt{config-vps}}

Skrypt \texttt{config-vps} łączy się z serwerem przez SSH i instaluje
wymagane oprogramowanie: powłokę \texttt{zsh}, agenta AI \texttt{gemini-cli}
oraz klienta \texttt{zerotier-one}. Tworzy katalog \texttt{~/bin} z programem
\texttt{ipb} (skrót do \texttt{ip -br}) i dołącza VPS do sieci ADM.
Reguły NSG w Azure (\texttt{LabSK-NSG}) przepuszczają ruch SSH, ICMP
oraz port ZeroTier UDP 9993.

\begin{lstlisting}
drapich@s3 ~ % config-vps
200 join OK
drapich@s3 ~ % ssh vps "echo $SHELL"
/usr/bin/zsh
drapich@s3 ~ % ssh vps "ipb a"
lo         UNKNOWN  127.0.0.1/8 ...
eth0       UP       10.0.0.4/24 ...
ztu7tkwebg UNKNOWN  10.80.108.71/24 ...
drapich@s3 ~ % ssh vps 'pacman -Q gemini-cli'
gemini-cli 1:0.45.2-1
\end{lstlisting}

\subsection{Sieć VPN/GAN -- ZeroTier ADM}

Sieć prywatna ADM została utworzona w panelu ZeroTier
(\texttt{my.zerotier.com}). Parametry sieci: NetID \texttt{8d1c312afa9a3adc},
zakres adresów \texttt{10.80.108.0/24}. Do sieci podłączono trzy urządzenia:
serwer VPS (\texttt{vps.adm}), laptop (\texttt{lap.adm}) oraz stację s3
(\texttt{sx.adm}). Każde urządzenie otrzymało stały adres IP w panelu
(Managed IPs) i zostało autoryzowane (Auth).

\begin{figure}[H]
  \centering
  \includesvg[width=0.58\textwidth]{ADM.svg}
  \caption{Schemat sieci ADM -- połączenia publiczne (niebieskie)
  i prywatne ZeroTier (zielone)}
\end{figure}

\subsection{Podłączenie stacji s3 do ADM}

Na stacji s3 zainstalowano klienta ZeroTier i dołączono ją do sieci
poleceniem \texttt{zerotier -n 8d1c312afa9a3adc add .}.
W panelu ZeroTier przypisano adres \texttt{10.80.108.19}.
Następnie dodano mapowania nazw \texttt{*.adm} w pliku \texttt{/etc/hosts}.

\begin{lstlisting}
drapich@s3 ~ % sudo pacman -S --needed zerotier-one
drapich@s3 ~ % zerotier -n 8d1c312afa9a3adc add .
drapich@s3 ~ % sudo zerotier-cli listnetworks
200 listnetworks 8d1c312afa9a3adc ADM ... OK ... 10.80.108.19/24
drapich@s3 ~ % ping -c1 vps.adm
64 bytes from 10.80.108.71: icmp_seq=1 ttl=64 time=18.2 ms
drapich@s3 ~ % ping -c1 lap.adm
64 bytes from 10.80.108.105: icmp_seq=1 ttl=64 time=22.1 ms
\end{lstlisting}

Połączenie publiczne (stacja s3 $\rightarrow$ VPS przez Internet):

\begin{lstlisting}
drapich@s3 ~ % ping -c1 vps
64 bytes from vps (74.248.150.53): icmp_seq=1 ttl=46 time=12.7 ms
drapich@s3 ~ % ssh vps hostname
arch
\end{lstlisting}

\section{Konfiguracja serwisu WWW}

\subsection{Nginx -- \texttt{config-www}}

Wybrany serwis: WWW (nginx). Skrypt \texttt{config-www} instaluje pakiet
\texttt{nginx} na serwerze VPS i włącza usługę w systemd.
Port TCP 80 jest otwarty w grupie bezpieczeństwa \texttt{LabSK-NSG}
(reguła \texttt{AllowHTTPInbound}). Serwer odpowiada domyślną stroną
\textit{Welcome to nginx!} zarówno po publicznym adresie \texttt{vps},
jak i po adresie ADM \texttt{vps.adm}.

\begin{lstlisting}
drapich@s3 ~ % config-www
Created symlink ... nginx.service
drapich@s3 ~ % ssh vps 'systemctl is-active nginx'
active
drapich@s3 ~ % curl -I http://vps
HTTP/1.1 200 OK
Server: nginx/1.30.2
drapich@s3 ~ % curl http://vps | head -3
<!DOCTYPE html>
<html>
<head>
\end{lstlisting}

\end{document}
